\section*{Broader Impacts and Limitations}

\textbf{Limitations.} While \textbf{RLearner-LLM} effectively mitigates stylistic hallucinations, several limitations warrant acknowledgement. First, the reliance on an external NLI model (DeBERTa-v3) bounded by its own training distribution presents a bottleneck: if the NLI model fails to grasp highly specialized domain jargon (e.g., advanced clinical oncology), the reward signal will inaccurately penalize valid reasoning. Second, NLI improvements under Hybrid-DPO are architecture-dependent: while LLaMA-2-13B shows consistent and large gains across all five domains, LLaMA-3-8B and Qwen3-8B exhibit mixed NLI results (with some domains showing regression relative to SFT), suggesting the alignment tax is most severe---and thus most addressable---when the base SFT model has limited initial logical grounding. Third, each domain test set comprises 100 held-out questions; future work should validate at larger scale with statistical significance tests. Finally, our reward scalarization (equal weighting of logic and fluency) may not represent the Pareto-optimal frontier for all educational contexts; younger learners may benefit from a heavier weighting toward fluency and narrative engagement.

\textbf{Broader Impacts.} The automated generation of educational explanations carries significant societal implications. On the positive side, \textbf{RLearner-LLM} democratizes access to high-fidelity, logically rigorous tutoring systems by allowing smaller, open-weight models (like LLaMA-2 13B) to achieve expert-level groundedness comparable to proprietary models. However, we must acknowledge the risk of over-reliance; while the alignment tax is reduced, LLMs inherently act as next-token predictors and remain susceptible to novel hallucination vectors. Deploying such systems in safety-critical domains like medical or legal examinations should always involve a human-in-the-loop validation protocol rather than autonomous grading.
